\begin{alist}
\item Έχουμε:
\[\log_2 8 = \log_2 2^3 = 3, \quad \log_2 \sqrt{2} = \log_2 2^{\frac{1}{2}} = \frac{1}{2} \quad \text{και} \quad \log_2 1 = 0.\]
Οπότε
\[\log_2 8 + 2\log_2 \sqrt{2} - \log_2 1 = 3 + 2 \cdot \frac{1}{2} - 0 = 3 + 1 = 4.\]

\item Ισχύει $\upsilon = P(\rho)$ με $\rho = 2$. \\
Είναι
\begin{align*}
\upsilon &= P(2) = (\log_2 8) \cdot 2^3 + 4(\log_2 \sqrt{2}) \cdot 2^2 - (4\log_2 1) \cdot 2 + 1990 \\
&= 8\log_2 8 + 16\log_2 \sqrt{2} - 8\log_2 1 + 1990 \\
&= 8(\log_2 8 + 2\log_2 \sqrt{2} - \log_2 1) + 1990 \\
&\overset{(\alpha)}{=} 8 \cdot 4 + 1990 \\
&= 2022
\end{align*}
Τελικά, το υπόλοιπο της διαίρεσης $P(x) : (x-2)$ είναι ίσο με $2022$.
\end{alist}
